\section{Related Work}

The closest prior line is spectral key-space steering. SEKA learns low-rank projections from synthetic question-answer pairs and injects them into attention keys during inference; AdaSEKA extends that idea with a small expert ensemble and an input-dependent expert selector \citep{li2026seka}. Those methods are designed for prompt highlighting and single-target emphasis. The operator itself is stationary: once a projection has been selected, the same amplified direction is applied throughout decoding. That design is natural when the objective is to preserve or recover one dominant target attribute.

Our paper differs first in the \emph{task}. We study multi-tool selection, where the model must emit two tools in sequence on the same prompt. This changes the failure mode. A method that keeps amplifying the same facet can look reasonable on a single-selection benchmark and still be mismatched for sequential coverage, because the second decoding step should become less like the first one, not more. For this reason we do not frame the paper as ``another SEKA variant.'' We frame it as an investigation of when the stationary key-side recipe stops matching the structure of the task.

The paper also differs in the \emph{intervention site}. The evaluated positive result comes from a query-side operator, not from key amplification. The operator is still low-rank and still uses an ontology-derived basis, so it is close enough to SEKA-style editing to make the comparison meaningful. At the same time, moving from keys to queries changes the sign semantics: a negative query-space coefficient suppresses already dominant ontology-aligned components and behaves like a coarse coverage prior. This is the core contrast we want the reader to keep in mind when interpreting the experiments.

More broadly, the paper belongs to post-hoc activation steering rather than training-time routing. The intervention is applied at inference time through attention-space hooks, with no additional finetuning and no external selector network. That choice is deliberate. It lets the experiment isolate whether the geometry of the ontology basis and the location of the intervention already determine the sign of the effect, before introducing history tracking or learned controllers.

KV-cache compression is orthogonal to the present manuscript. Methods such as KIVI study low-bit cache representations and memory-quality tradeoffs rather than tool-selection dynamics \citep{liu2024kivi}. Earlier drafts in this repository explored that direction, but the current paper is intentionally scoped away from compression. The evaluated models are instruction-tuned Qwen, Llama, and Mistral variants \citep{qwen2024qwen25,grattafiori2024llama3,jiang2023mistral}, and the main benchmark is MetaTool Subtask4 because it directly exposes the sequential-coverage problem that stationary single-facet steering is expected to mishandle.
